\section{Related Work}
\label{sec:related_work}

\paragraph{Reinforcement learning with verifiable rewards.}
RLHF~\citep{ouyang2022training} established the paradigm of aligning
language models via reward signals derived from human preferences.
DeepSeek-R1~\citep{deepseekr1} and
DeepSeekMath~\citep{shao2024deepseekmath} extended this to
\emph{verifiable} rewards (RLVR), using programmatic correctness checks
for mathematical reasoning. GRPO~\citep{shao2024deepseekmath} computes
group-relative advantages without a critic network, offering a
lightweight alternative to PPO. These approaches rely exclusively on
outcome-level signals and do not supervise the intermediate reasoning
process. TopoPRM complements RLVR by adding process-level topology
rewards that remain fully verifiable and code-computed.

\paragraph{Process reward models.}
\citet{lightman2023lets} introduced process reward models (PRMs) that
assign scalar rewards to each reasoning step, trained via human
annotations of step-level correctness. Math-Shepherd~\citep{wang2024math}
automates PRM data collection using Monte Carlo estimation.
Proof-RM~\citep{proofRM2026} designs scalable rewards specifically for
mathematical proofs. While PRMs provide step-level supervision, they
require either expensive human labels or learned verifiers that may
themselves be unreliable. In contrast, TopoPRM derives process rewards
from structural properties of the reasoning graph, requiring no learned
components and introducing no additional model parameters.

\paragraph{Graph-structured reasoning.}
Graph-of-Thought~\citep{besta2024graph} organises reasoning into
graph-like structures via prompt engineering.
GRP~\citep{grp2025} proposes a graph reasoning paradigm that trains
models to produce graph-structured outputs with cognitive labels. While
GRP focuses on \emph{generating} graph-structured solutions with
specialised labels, TopoPRM uses graph topology as a \emph{reward
signal} for standard autoregressive generation, requiring no
architectural modifications or specialised output formats.

\paragraph{Reward hacking and advantage estimation.}
Reward hacking---where models exploit auxiliary rewards without
improving task performance---is a well-known challenge in
RL~\citep{Denison2024SycophancyTS}. Recent work on GRPO variants
addresses related issues: DAPO introduces dynamic
sampling~\citep{dapo2026}, while difficulty-aware
GRPO~\citep{difficultyGRPO2026} adjusts training based on problem
difficulty. Our SCAE method is orthogonal to these approaches,
specifically addressing the interaction between outcome and process
rewards through stratified advantage clipping.

\paragraph{Reasoning compression.}
On-policy self-distillation~\citep{reasoning_compression_2025} and
reverse-KL distillation~\citep{hinton2015distilling} have been explored
for compressing verbose chain-of-thought reasoning into shorter outputs.
Our distillation stage transfers both accuracy and structural conciseness
from the TopoPRM-trained teacher to smaller student models.
